\section{Exterior interactions and Rodgers--Tao renormalization}
\label{sec:RT}

\subsection{What the real-zero argument controls}
The factorization in Section~\ref{sec:double} removes the singular
mutual interaction of one pair. In the real Laguerre--P\'olya setting,
\(-\kappa\) is an exterior inverse-square sum, with the possible
Gaussian term in \eqref{eq:LP-kappa}. This connects the critical
coefficients to the interaction quantities in Rodgers and Tao's
proof \cite{RodgersTao2020}. Their theorem proves \(\Lambda\ge0\)
by contradiction from \(\Lambda<0\): real, simple zero dynamics
then exists on an interval reaching the zeta slice \(t=0\).
Here their parameter \(t\) is our Newman parameter \(\lambda\),
and has the opposite orientation to \(\tau\).

Their finite-block localization and renormalized energies control
deviations from classical spacing over growing blocks. The
definitions and estimates used here are those of corrected arXiv
version~5, Sections~4 and~7; HTML theorem numbers differ from the
journal/PDF numbering. We use the method and its scope, not an
unstated extension of its estimates to \(\Lambda>0\).

\subsection{Renormalization retains the collision pole}
For a positive actual separation \(g\) and reference separation
\(d>0\), the pair interaction is renormalized using
\begin{equation}\label{eq:RT-renorm}
 V(r)=r^{-2}-1+2(r-1),\qquad
 \widetilde E(g,d)=d^{-2}V(g/d).
\end{equation}
Since \(V(1)=V'(1)=0\) and \(V''(r)=6r^{-4}>0\), this is
nonnegative for \(r>0\). It subtracts the constant and linear
reference contribution, not the singular interaction. Indeed
\begin{equation}\label{eq:RT-pole}
 \widetilde E(g,d)=g^{-2}-3d^{-2}+2g d^{-3}
                =g^{-2}+O(1)\quad(g\downarrow0)
\end{equation}
at fixed \(d\). Ordered-pair sums count this term twice.
Neither \(\widetilde E\) nor \(-\kappa\) is the field energy
\(\inner\Psi{K\Psi}\) or the local quotient
\(T_{\mathrm{loc}}/M_{\mathrm{loc}}\).

The local deflation \(H=P U\) and the reference subtraction
\eqref{eq:RT-renorm} perform different jobs. The former gives a
regular local coefficient through a collision; the latter makes
a large real-zero interaction comparable with an equilibrium
configuration. An integrated estimate over growing blocks does
not by itself bound one prescribed exterior coefficient or
prevent an isolated finite-height event.

This leaves two distinct questions. One is whether the exact theta
source gives effective bounds on \(\ell,\kappa\) in a declared
cluster region, using \eqref{eq:moment-ratios} or
\eqref{eq:contour-exterior}. The other is whether a statistical
argument can be formulated when the zeros are not all real.
The latter requires retaining the transverse zero coordinates;
it is taken up in Section~\ref{sec:complex-statistics}. No extension of the Rodgers--Tao
relaxation theorem is presumed in making that definition.
